\documentclass[12pt, letterpaper]{article}
\usepackage[utf8]{inputenc}
\usepackage[spanish]{babel}
\usepackage{graphicx}
\usepackage{amsmath}
\usepackage{amsfonts}
\usepackage{amssymb}
\usepackage{geometry}
\usepackage{setspace}
\usepackage{tocloft}
\usepackage{parskip}
\usepackage{hyperref}
\usepackage{fancyhdr}
\usepackage{array} % Paquete para mejorar las tablas
\usepackage{listings} % Paquete para bloques de código
\usepackage{xcolor}   % Paquete para colores en el código
\usepackage{float}    % Paquete para mejor posicionamiento de figuras [H]
\usepackage[style=apa, backend=biber]{biblatex}
\addbibresource{references.bib} % Archivo .bib para las referencias

\geometry{
    letterpaper,
    total={170mm,257mm},
    left=30mm,
    right=30mm,
    top=25mm,
    bottom=25mm,
}
\renewcommand{\cftsecleader}{\cftdotfill{\cftdotsep}}
\setlength{\cftbeforesecskip}{1em}

%=========== ESTILO PARA EL CÓDIGO ===========
\definecolor{codegreen}{rgb}{0,0.6,0}
\definecolor{codegray}{rgb}{0.5,0.5,0.5}
\definecolor{codepurple}{rgb}{0.58,0,0.82}
\definecolor{backcolour}{rgb}{0.95,0.95,0.92}

\lstdefinestyle{sqlstyle}{
    backgroundcolor=\color{backcolour},   
    commentstyle=\color{codegreen},
    keywordstyle=\color{magenta},
    numberstyle=\tiny\color{codegray},
    stringstyle=\color{codepurple},
    basicstyle=\footnotesize\ttfamily,
    breakatwhitespace=false,         
    breaklines=true,                 
    captionpos=b,                    
    keepspaces=true,                 
    numbers=left,                    
    numbersep=5pt,                  
    showspaces=false,                
    showstringspaces=false,
    showtabs=false,                  
    tabsize=2,
    language=SQL
}

\lstdefinestyle{dbmlstyle}{
    backgroundcolor=\color{backcolour},
    commentstyle=\color{codegreen},
    keywordstyle=\color{blue},
    stringstyle=\color{codepurple},
    basicstyle=\footnotesize\ttfamily,
    breaklines=true,
    captionpos=b,
    keepspaces=true,
    numbers=left,
    numbersep=5pt,
    showstringspaces=false,
    morekeywords={Table, integer, primary, key, increment, varchar, unique, not, null, ref, numeric}
}


\pagestyle{fancy}
\fancyhf{}
\renewcommand{\headrulewidth}{0pt}
\cfoot{\thepage}

\begin{document}

\begin{titlepage}
    \centering
    \vspace*{1cm}
    {\Huge \bfseries Software de implementación de funcionalidades y aplicaciones de un modelo matemático para el análisis y optimización de mezclas de concreto estructurales.\par}
    \vspace{2.5cm}
    {\Large \bfseries Modalidad de trabajo de grado\par}
    \vspace{0.5cm}
    {\Large Auxiliar de Semillero de Investigación\par}
    \vspace{2.5cm}
    {\Large \bfseries Juan Diego Pedroza Sierra\par}
    \vspace{4cm}
    {\Large Universidad Cooperativa de Colombia\par}
    {\Large Facultad de Ingeniería\par}
    {\Large Programa de Ingeniería Sistemas\par}
    {\Large Villavicencio, Meta\par}
    {\Large 2023\par}
\end{titlepage}



\newpage
\thispagestyle{empty}
\begin{center}
    \vspace*{2cm}
    {\large \bfseries Autoridades académicas Universidad Cooperativa de Colombia\par}
    \vspace{2cm}
\end{center}
\begin{flushleft}
    \large
    Dra. Maritza Rondón Rangel\\
    \textit{Rectora Nacional}\\[1cm]
    Dr. César Augusto Pérez Londoño\\
    \textit{Director de la Sede}\\[1cm]
    Dr. Henry Emiro Vergara Bobadilla\\
    \textit{Subdirector Académico de la Sede}\\[1cm]
    Dra. Ruth Edith Muñoz\\
    \textit{Directora Administrativa}\\[1cm]
    Ing. Raúl Alarcón Bermúdez\\
    \textit{Decano Facultad de Ingeniería Civil}\\[1cm]
    Ing. María Lucrecia Ramírez Suárez\\
    \textit{Jefe programa de Ingenierías}\\[1cm]
    Ing. Pedro Alexander Gutiérrez Aguilera\\
    \textit{Coordinador de Investigación del Programa de Ingeniería Civil}\\[1cm]
\end{flushleft}

\newpage
\thispagestyle{empty}
\begin{center}
    \vspace*{3cm}
    {\large \bfseries Advertencia\par}
    \vspace{2cm}
    \large
    La Universidad Cooperativa de Colombia, sede Villavicencio, no se hace responsable por los conceptos emitidos por los autores.
\end{center}

\newpage
\thispagestyle{empty}
\begin{center}
    \vspace*{3cm}
    {\large \bfseries Nota de aceptación\par}
    \vspace{2cm}
    \rule{15cm}{0.4pt}
    \rule{15cm}{0.4pt}
    \rule{15cm}{0.4pt}
    \rule{15cm}{0.4pt}
    \vspace{4cm}
\end{center}
\begin{flushleft}
    \large
    \begin{tabular}{p{7cm}p{7cm}}
        \rule{7cm}{0.4pt} & \rule{7cm}{0.4pt} \\
        Firma del presidente del jurado & Firma del jurado \\
    \end{tabular}
    \vspace{2cm}
    \begin{center}
        \rule{7cm}{0.4pt}\\
        Firma del jurado
    \end{center}
\end{flushleft}

\newpage
\thispagestyle{empty}
\begin{flushright}
    \vspace*{3cm}
    {\large \bfseries Dedicatoria\par}
    \vspace{2cm}
    \textit{Esta página es opcional….\\ 
    Mencione, en pocas líneas, las personas a quienes dedica su trabajo. Ubique el texto alineado a la derecha sin justificar.}
\end{flushright}

\newpage
\thispagestyle{empty}
\begin{center}
    \vspace*{3cm}
    {\large \bfseries Agradecimientos\par}
    \vspace{2cm}
    \textit{Esta página es opcional. Mencione aquellas personas, instituciones u organizaciones que hicieron posible el desarrollo del trabajo de grado.}
\end{center}

\newpage
\pagenumbering{roman}
\setcounter{page}{1}
\section{Resumen}

En el departamento del Meta (Colombia), las condiciones ambientales aceleran la degradación del concreto y elevan los costos de mantenimiento. Ante la ausencia de herramientas locales para apoyar la dosificación y selección de aditivos, este proyecto desarrolla DIAPCE: una aplicación multiplataforma en Flutter que asiste la toma de decisiones para mezclas de concreto estructural y está orientada al trabajo en los laboratorios de materiales y a los semilleros de investigación de la Universidad Cooperativa de Colombia. A partir de un conjunto de datos regional y entradas como resistencia objetivo, temperatura, humedad y condiciones de exposición, el sistema sugiere una dosificación base y clasifica el uso de aditivos (P0–PP3), con proyecciones a 7, 14 y 28 días, en concordancia con ACI 211.1 y normas relacionadas. La solución funciona sin conexión mediante SQLite y prioriza la trazabilidad de los ensayos y ejercicios académicos (registro de insumos, resultados y versión del modelo), facilitando prácticas reproducibles y comparables en el aula y el laboratorio. DIAPCE busca mejorar desempeño y durabilidad y, potencialmente, contribuir a la sostenibilidad al fomentar un uso más eficiente de materiales.

\vspace{0.5em}
\noindent\textbf{Palabras clave:} concreto estructural, dosificación, aditivos (P0–PP3), ACI 211.1, Flutter, SQLite, Meta (Colombia), 7/14/28 días, laboratorios, semilleros de investigación, docencia.

\newpage
\tableofcontents
\newpage
\listoftables
\addcontentsline{toc}{section}{Lista de tablas}
\newpage
\listoffigures
\addcontentsline{toc}{section}{Lista de figuras}

\newpage
\pagenumbering{arabic}
\setcounter{page}{1}


% docs/planteamiento.tex
\section{Planteamiento del problema}

En el departamento del Meta (Colombia), las condiciones ambientales agresivas —altas temperaturas, humedad elevada y particularidades de los suelos— inciden de manera directa en el desempeño y la durabilidad del concreto empleado en obras civiles. A pesar de ello, existe una carencia de herramientas digitales locales que apoyen, de forma didáctica y con trazabilidad, la dosificación y la selección de aditivos según el contexto regional. Esta brecha se hace evidente en los laboratorios de materiales y en los semilleros de investigación de la Universidad Cooperativa de Colombia (UCC), donde la replicabilidad de los ensayos y la sistematización de resultados son esenciales para la formación y la transferencia de conocimiento.

La ausencia de software especializado para apoyar el diseño y la verificación de mezclas adaptadas a condiciones locales puede derivar en decisiones inconsistentes, menor durabilidad y mayores costos de mantenimiento. De acuerdo con lineamientos nacionales, es prioritario impulsar la innovación tecnológica en el sector construcción para elevar la calidad y sostenibilidad de la infraestructura (Departamento Nacional de Planeación, 2018). A su vez, gremios y estudios académicos resaltan la necesidad de incorporar herramientas digitales que integren particularidades ambientales y estándares técnicos en el diseño de mezclas para optimizar recursos y asegurar la vida útil de las estructuras (CCI, 2020; González et al., 2019).

En este contexto, se identifica la necesidad de una solución que permita: (i) registrar de manera estructurada las condiciones de diseño (resistencia objetivo, temperatura, humedad, exposición), (ii) asistir la dosificación base y la clasificación del uso de aditivos (P0–PP3) con proyecciones a 7, 14 y 28 días, en concordancia con ACI 211.1 y normas relacionadas, y (iii) garantizar la trazabilidad de los ensayos y ejercicios académicos. DIAPCE responde a esta necesidad mediante una aplicación multiplataforma desarrollada en Flutter que opera sin conexión (SQLite), orientada a laboratorios y semilleros de la UCC para apoyar decisiones basadas en datos, estandarizar prácticas y facilitar la comparación reproducible de resultados. La herramienta no sustituye el criterio del ingeniero ni el marco normativo, sino que lo complementa con evidencia sistematizada y accesible en el contexto regional del Meta.


\subsection*{Pregunta de investigación}
¿Cómo por medio de la implementación de un software se pueden realizar cálculos de optimización de mezclas cementantes a partir de un modelo matemático que use propiedades físicas, mecánicas y térmicas de concreto estructural?

\section{Objetivos}
\subsection{Objetivo general}
Desarrollar un software que implemente un modelo matemático para el ajuste de proporciones de materiales cementantes en mezclas de concreto estructural.

\subsection{Objetivos específicos}
\begin{enumerate}
    \item Implementar una base de datos que permita el almacenamiento de características físicas, mecánicas y/o ambientales de mezclas de concreto estructural.
    \item Desarrollar un software que implemente un modelo matemático para las mezclas de concreto estructural con una arquitectura cliente servidor.
    \item Realizar pruebas y análisis de desempeño del programa referente a criterios de calidad.
\end{enumerate}

\section{Introducción}
La industria de la construcción en el departamento del Meta, Colombia, se enfrenta a retos constantes derivados de las condiciones ambientales de la región, las cuales exigen materiales de alta durabilidad y resistencia. El concreto, siendo el material estructural por excelencia, requiere de un diseño de mezcla preciso que se adapte a estas particularidades para garantizar la longevidad y seguridad de las infraestructuras. Actualmente, existe una brecha en la disponibilidad de herramientas tecnológicas que permitan una optimización eficiente de estas mezclas.

Este trabajo de grado aborda dicha problemática mediante el desarrollo de una aplicación de software denominada \textbf{DIAPCE (Diseño y Análisis de Proporciones de Concreto Estructural)}. El objetivo principal de DIAPCE es facilitar a profesionales y estudiantes del sector (ingenieros, constructores) la creación y optimización de diseños de mezcla de concreto estructural, especialmente adaptados a las condiciones del Meta. La aplicación permite registrar, analizar y optimizar mezclas, almacenando información clave sobre materiales, proporciones, propiedades físicas, mecánicas y térmicas.

El documento se estructura de la siguiente manera: se inicia con el planteamiento del problema y la justificación. Posteriormente, se presenta el marco teórico que sustenta el desarrollo, abarcando modelos matemáticos para la optimización de concreto y las tecnologías seleccionadas. Se detalla la metodología, incluyendo la arquitectura del sistema y el enfoque Feature-Driven Development (FDD). Finalmente, se exponen los resultados esperados, el plan de trabajo y el presupuesto asociado al proyecto.

\section{Marco teórico}
\subsection{Diseño de mezclas y ley de Abrams}
El diseño de mezclas de concreto se apoya históricamente en relaciones empíricas entre la resistencia a compresión y la relación agua/cemento (a/c). La ley de Abrams propone una relación inversa entre ambas magnitudes, que en forma log-lineal puede expresarse como:
\[
f'_c = K \cdot \left(\frac{w}{c}\right)^{-n}
\quad\text{o}\quad
\log f'_c = \log K - n \cdot \log\!\left(\frac{w}{c}\right),
\]
donde $f'_c$ es la resistencia característica, $K$ y $n$ son parámetros dependientes de materiales y condiciones de curado. Métodos normativos (p. ej., ACI 211.1) emplean esta y otras relaciones para estimar dosificaciones iniciales que luego se ajustan con ensayos. En DIAPCE, la relación a/c es un nodo clave del flujo de decisión y se restringe a valores presentes en el conjunto experimental, reduciendo la incertidumbre del usuario al seleccionar solo combinaciones observadas.

\subsection{Efecto de condiciones ambientales (temperatura y humedad)}
La cinética de hidratación del cemento y, por tanto, el desarrollo de la resistencia, dependen de la temperatura y el estado higrotérmico. Temperaturas más altas aceleran la hidratación inicial pero pueden reducir resistencias finales si hay pérdida de humedad; ambientes más secos favorecen gradientes de humedad y retracción. Por esto, DIAPCE incorpora explícitamente:
\begin{itemize}
    \item \textbf{Temperatura ambiente (°C)}: primer filtro del flujo; acota el subconjunto de registros.
    \item \textbf{Humedad relativa (\%)}: segundo filtro; condiciona las opciones de relación a/c.
\end{itemize}
Esta cascada emula un \textit{árbol de decisión} que condiciona las alternativas posteriores a partir de parámetros ambientales, manteniendo la coherencia con los datos observados.

\subsection{Aditivos químicos y códigos P0–PP3}
Los aditivos modifican trabajabilidad, tiempo de fraguado y desarrollo de resistencia (p. ej., reductores de agua, superplastificantes, acelerantes, retardantes; ASTM C494). En el dataset se normalizan bajo códigos \texttt{P0}–\texttt{PP3}, lo que permite:
\begin{itemize}
    \item \textbf{Selección guiada}: se muestran solo aditivos válidos para la combinación temperatura–humedad–a/c.
    \item \textbf{Trazabilidad}: el código simplifica el JOIN con tablas técnicas (\texttt{tipos\_aditivo}, \texttt{productos}, \texttt{aditivos}).
\end{itemize}

\subsection{Evolución temporal de la resistencia (7, 14 y 28 días)}
La resistencia se reporta típicamente a 7, 14 y 28 días (ASTM C39). A 28 días se toma como referencia de diseño. DIAPCE estima las resistencias $\{f_7, f_{14}, f_{28}\}$ como promedios empíricos del subconjunto filtrado por condiciones ambientales, relación a/c y aditivo:
\[
\hat f_{t} \;=\; \operatorname{avg}\{\, f_{t,i}\,:\, (T_i,H_i,\mathrm{a/c}_i,\mathrm{aditivo}_i)\in \text{selección}\,\},\quad t\in\{7,14,28\}.
\]
Este enfoque \textit{data-driven} privilegia la reproducibilidad y la coherencia con el comportamiento observado en laboratorio del conjunto local (1874 registros).

\subsection{Modelado empírico y enfoques de aprendizaje}
La literatura reciente explora desde regresiones multivariadas hasta redes neuronales para predecir $f'_c$ incorporando variables como a/c, tipo de cemento, granulometría, aditivos y condiciones de curado. DIAPCE adopta en esta versión un \textbf{modelo empírico de agregación} (promedios condicionados) por su:
\begin{itemize}
    \item \textbf{Transparencia}: resultados explicables y trazables a registros concretos.
    \item \textbf{Robustez} ante datos ruidosos cuando se hace estratificación adecuada.
    \item \textbf{Compatibilidad} con decisiones en cascada y validaciones de rango (resistencia objetivo 24–57 MPa con decimales).
\end{itemize}
La arquitectura y el esquema de datos dejan abierta la integración futura de modelos ML (p. ej., regresión regularizada o \textit{tree ensembles}) entrenados sobre las mismas vistas SQL.

\subsection{Sistemas de soporte a decisiones y consultas relacionales}
DIAPCE actúa como un \textit{asistente de decisión} local. Su núcleo es una base de datos SQLite con:
\begin{itemize}
    \item \textbf{Normalización y llaves foráneas} para garantizar integridad (proyectos, mezclas, materiales, aditivos, resultados de concreto).
    \item \textbf{Índices} (p. ej., \texttt{idx\_busqueda\_resistencia}) para acelerar las consultas condicionales.
    \item \textbf{Vistas/consultas} que implementan el árbol de decisión: 
    temperatura $\rightarrow$ humedad $\rightarrow$ relación a/c $\rightarrow$ aditivo, 
    y las funciones de predicción $\{\hat f_7,\hat f_{14},\hat f_{28}\}$.
\end{itemize}
Este patrón \textit{relational-first} asegura reproducibilidad de resultados y portabilidad sin dependencia de red.

\subsection{Arquitectura móvil y visualización}
La aplicación se implementa en Flutter (Dart), con almacenamiento local (\texttt{sqflite}) y visualizaciones con \texttt{fl\_chart}:
\begin{itemize}
    \item \textbf{UI guiada}: formularios con validadores que restringen a intervalos y opciones presentes en el dataset.
    \item \textbf{Gráficos}: pastel para composición de mezcla y línea para la evolución de resistencia, apoyando interpretación técnica.
    \item \textbf{Proyectos por usuario}: aislamiento lógico que facilita trazabilidad y organización de diseños.
\end{itemize}

\subsection{Síntesis conceptual aplicada en DIAPCE}
Integrando los elementos anteriores, el marco teórico de DIAPCE se puede resumir en tres pilares:
\begin{enumerate}
    \item \textbf{Fundamento de materiales}: relación a/c (ley de Abrams), influencia higrotérmica, y efecto de aditivos en la hidratación y trabajabilidad.
    \item \textbf{Evidencia experimental local}: predicciones y opciones guiadas extraídas de un dataset real de 1874 registros, con estratificación por condiciones.
    \item \textbf{Soporte computacional}: consultas relacionales optimizadas, decisión en cascada y visualizaciones que transforman datos en evidencia accionable.
\end{enumerate}

\subsection{Trabajos relacionados y estándares de referencia}
Como línea base se consideran los métodos ACI para diseño de mezclas (ACI 211.1) y normas de ensayo (ASTM C39 para compresión, ASTM C494 para aditivos), junto con literatura de optimización y predicción de resistencia (p. ej., relaciones empíricas y regresión multivariable). DIAPCE se posiciona como un \textit{asistente data-driven} local, complementario a softwares comerciales, con énfasis en adaptabilidad regional y trazabilidad de decisiones.




% docs/justificacion.tex
\section{Justificación}

La construcción en regiones como el Meta enfrenta retos estructurales: condiciones ambientales agresivas, asimetrías en la disponibilidad de datos locales y brechas de estandarización en prácticas de laboratorio y formación. Estos factores dificultan decisiones consistentes de dosificación y selección de aditivos y, a largo plazo, impactan la durabilidad y los costos de mantenimiento. En el entorno académico, los laboratorios de materiales y los semilleros de investigación de la UCC requieren herramientas que faciliten el registro disciplinado, la comparación reproducible de resultados y la transferencia de conocimiento.

En este contexto, la implementación de una herramienta como DIAPCE es pertinente por tres razones principales: (i) alinea la práctica con la normativa (p.\,ej., ACI 211.1 y normas relacionadas) y estandariza el flujo de decisión mediante entradas explícitas (resistencia objetivo, temperatura, humedad, exposición); (ii) promueve la trazabilidad al almacenar insumos, resultados y versión del modelo de manera local (SQLite), favoreciendo la replicabilidad de ensayos, la evaluación docente y la auditoría académica; y (iii) opera sin conexión, lo que resulta clave en escenarios con limitaciones de conectividad, sin fricción de despliegue en laboratorios.

La elección tecnológica responde a criterios de portabilidad, facilidad de adopción y costo total: Flutter permite un desarrollo multiplataforma desde un único código base, facilitando el uso en equipos heterogéneos (dispositivos móviles y de escritorio), mientras que SQLite habilita una base de datos embebida, liviana y suficiente para la gestión de proyectos y mezclas en contextos educativos. Este diseño centrado en cliente es coherente con el alcance actual del proyecto y minimiza dependencias externas; adicionalmente, deja abierta la posibilidad de una evolución futura hacia servicios remotos (p.\,ej., sincronización, analítica o modelos avanzados) sin comprometer el uso inmediato en laboratorio.

Desde una perspectiva de impacto, DIAPCE contribuye a la formación de competencias en diseño de mezcla, análisis de resultados y cultura de registro, apoyando a estudiantes y docentes en la toma de decisiones basadas en datos. Potencialmente, el uso sistemático de la herramienta puede incentivar un empleo más eficiente de materiales y la selección adecuada de aditivos, con beneficios asociados en desempeño y sostenibilidad. Cabe enfatizar que DIAPCE complementa —no reemplaza— el criterio profesional ni el cumplimiento normativo: su propósito es ofrecer evidencia organizada y accesible que fortalezca la práctica académica y técnica en el contexto regional del Meta.


\section{Metodología}
Como metodología de trabajo optamos por usar FDD (Feature-Driven Development) debido a su enfoque ágil y centrado en el cliente. FDD facilita la división del trabajo en pequeñas funciones, permitiendo una detección temprana de errores y una rápida adaptación a los cambios (Lucidchart, 2024). El tipo de investigación es mixto, con un enfoque en el desarrollo de software.

\begin{figure}[H]
    \centering  
    % NOTA: Asegúrate de tener el archivo 'diagrama_db.jpg' en la misma carpeta.
    \includegraphics[width=1\textwidth]{images/metodolgia.jpg}
    \caption{Diagrama metodologico del proyecto con la estructura de la documentación y desarrollo.}
    \label{fig:diagrama}
\end{figure}

\subsection{Planteamiento de la Solución: DIAPCE}
DIAPCE es una aplicación desarrollada en Flutter, diseñada como una herramienta de apoyo para el diseño de mezclas de concreto, fundamentada en datos experimentales locales.

\subsubsection*{Propósito y Flujo Principal}
El objetivo de DIAPCE es asistir a ingenieros y técnicos en la selección de mezclas de concreto, basando las recomendaciones en datos reales en lugar de estimaciones. La aplicación permite a los usuarios crear y gestionar proyectos, donde pueden definir condiciones controladas (temperatura, humedad, relación agua/cemento, aditivos) y obtener predicciones de resistencia a los 7, 14 y 28 días. Estas predicciones se calculan a partir de una base de datos local, poblada a través de la importación de un archivo CSV con resultados de ensayos experimentales.

El flujo de usuario es el siguiente:
\begin{enumerate}
    \item El usuario inicia la aplicación e ingresa a su perfil.
    \item Crea un nuevo proyecto, especificando nombre, fecha, tipo de obra y una imagen opcional.
    \item Define las propiedades técnicas deseadas mediante un sistema de selección en cascada: resistencia objetivo, temperatura, humedad, relación a/c y aditivo.
    \item La aplicación consulta la base de datos experimental, calcula los promedios de resistencia para las condiciones seleccionadas y presenta una predicción gráfica y numérica.
    \item El proyecto, junto con sus condiciones y predicciones, se guarda en la base de datos local SQLite, asociado al perfil del usuario.
\end{enumerate}

\subsection{Arquitectura y Tecnologías}
La aplicación DIAPCE se desarrolla bajo una arquitectura cliente-servidor, diseñada para ser escalable, segura y de fácil mantenimiento.
\begin{itemize}
    \item \textbf{Frontend (Cliente):} Desarrollado en \textbf{Flutter} con el lenguaje \textbf{Dart}. Esta elección permite crear una experiencia de usuario nativa y fluida en múltiples plataformas. Para la persistencia de datos locales (proyectos de usuario), se utiliza una base de datos \textbf{SQLite} a través del paquete `sqflite`.
    \item \textbf{Backend (Servidor):} Implementado en \textbf{Python}, utilizando un framework como \textbf{Flask} o \textbf{Django}. El backend expone una API RESTful que procesa la lógica de negocio, incluyendo los cálculos del modelo matemático y la gestión de la base de datos principal.
    \item \textbf{Base de Datos (Servidor):} Se utilizará \textbf{PostgreSQL}, donde residirán los datos experimentales consolidados y la información centralizada del sistema.
    \item \textbf{Comunicación:} La comunicación entre el cliente y el servidor se realiza mediante peticiones HTTP a la API RESTful, utilizando JSON como formato de intercambio de datos.
\end{itemize}

\subsection{Variables de investigación}
\begin{itemize}
    \item \textbf{Variable Independiente:} El modelo matemático y los parámetros de entrada.
    \item \textbf{Variable Dependiente:} Los resultados generados por el software (proporciones de mezcla) y las métricas de calidad.
    \item \textbf{Variables Intervinientes:} Errores en los datos de entrada, problemas de compatibilidad entre tecnologías.
    \item \textbf{Variables Controladas:} Formatos de entrada de datos, entorno de pruebas y criterios de desempeño definidos.
\end{itemize}

\subsection{Actividades}
El proyecto se divide en las siguientes fases:
% ... (Las fases se mantienen como estaban) ...

\subsection{Estudio de Caso: Implementación de la Base de Datos Experimental}
Para cumplir con el primer objetivo específico, se realizó un proceso detallado de diseño e implementación de la base de datos en PostgreSQL. Este proceso partió de un conjunto de datos brutos provenientes de ensayos de resistencia, los cuales presentaban desafíos significativos.

\subsubsection{Análisis del Problema y Normalización de Datos}
Los datos originales presentaban problemas de ambigüedad y redundancia, como el uso de un mismo código para productos diferentes y la repetición de nombres comerciales en miles de filas. Para solucionar esto, se aplicó un enfoque de normalización progresiva:
\begin{enumerate}
    \item \textbf{Primera Etapa (3 Tablas):} Se separaron los tipos de aditivos y sus códigos en tablas de catálogo para resolver la ambigüedad principal.
    \item \textbf{Refinamiento (4 Tablas):} Se identificó que los nombres comerciales de los productos aún generaban redundancia. Para alcanzar una mayor normalización, se creó una tabla dedicada (`Productos`), resultando en un diseño final de cuatro tablas que garantiza la integridad y escalabilidad de los datos.
\end{enumerate}
Antes de la importación, los datos fueron limpiados y estandarizados en formato CSV, asignando códigos únicos y mapeando valores para su correcta inserción.

\subsubsection{Diseño del Esquema Final y Diagrama ER}
El esquema final, altamente normalizado, se compone de cuatro tablas interrelacionadas. El siguiente diagrama (Figura \ref{fig:diagrama}) ilustra la estructura.

\begin{figure}[H]
    \centering  
    % NOTA: Asegúrate de tener el archivo 'diagrama_db.jpg' en la misma carpeta.
    \includegraphics[width=1\textwidth]{images/diagrama_db.png}
    \caption{Diagrama Entidad-Relación de la estructura final de 4 tablas.}
    \label{fig:diagrama}
\end{figure}

El código en DBML (Database Markup Language) utilizado para generar este diagrama se muestra a continuación, detallando las columnas, tipos de datos y relaciones.

\begin{lstlisting}[style=dbmlstyle, caption=Código en DBML para la herramienta dbdiagram.io.]
Table TiposAditivo {
  id integer [primary key, increment]
  nombre varchar(50) [unique, not null]
}

Table Productos {
  id integer [primary key, increment]
  nombre_producto varchar(100) [unique, not null]
  marca varchar(50)
}

Table Aditivos {
  id integer [primary key, increment]
  codigo varchar(10) [unique, not null]
  porcentaje_aplicado varchar(10) [not null]
  tipo_aditivo_id integer [not null, ref: > TiposAditivo.id]
  producto_id integer [not null, ref: > Productos.id]
}

Table ResultadosConcreto {
  id integer [primary key, increment]
  temperatura integer [not null]
  humedad integer [not null]
  relacion_ac numeric(3,2) [not null]
  edad_dias integer [not null]
  resistencia_mpa numeric(10,2) [not null]
  aditivo_id integer [not null, ref: > Aditivos.id]
}
\end{lstlisting}

\subsubsection{Implementación en PostgreSQL y Consultas de Análisis}
El esquema se implementó en PostgreSQL utilizando sentencias DDL (Data Definition Language) para crear la estructura de tablas.

\begin{lstlisting}[style=sqlstyle, caption=Script de creación de las 4 tablas en PostgreSQL.]
-- Crear tablas de catálogo
CREATE TABLE TiposAditivo (
    id SERIAL PRIMARY KEY,
    nombre VARCHAR(50) UNIQUE NOT NULL
);
CREATE TABLE Productos (
    id SERIAL PRIMARY KEY,
    nombre_producto VARCHAR(100) UNIQUE NOT NULL,
    marca VARCHAR(50)
);
-- Crear tabla de enlace de aditivos
CREATE TABLE Aditivos (
    id SERIAL PRIMARY KEY,
    codigo VARCHAR(10) UNIQUE NOT NULL,
    porcentaje_aplicado VARCHAR(10) NOT NULL,
    tipo_aditivo_id INT NOT NULL REFERENCES TiposAditivo(id),
    producto_id INT NOT NULL REFERENCES Productos(id)
);
-- Crear tabla principal de resultados
CREATE TABLE ResultadosConcreto (
    id SERIAL PRIMARY KEY,
    temperatura INT NOT NULL,
    humedad INT NOT NULL,
    relacion_ac DECIMAL(3, 2) NOT NULL,
    edad_dias INT NOT NULL,
    resistencia_mpa DECIMAL(10, 2) NOT NULL,
    aditivo_id INT NOT NULL REFERENCES Aditivos(id)
);
\end{lstlisting}

Una vez creadas las tablas, se poblaron con los datos maestros y, posteriormente, con los datos de los ensayos. Esta estructura normalizada permite realizar consultas de análisis complejas y eficientes. Por ejemplo, la siguiente consulta calcula la resistencia promedio para cada grupo experimental, redondeando el resultado a dos decimales.

\begin{lstlisting}[style=sqlstyle, caption=Cálculo de la resistencia promedio por grupo experimental.]
SELECT
    p.nombre_producto,
    a.codigo AS codigo_aditivo,
    r.edad_dias,
    -- Se redondea el promedio a 2 decimales
    ROUND(AVG(r.resistencia_mpa), 2) AS promedio_resistencia,
    COUNT(r.id) AS numero_de_muestras
FROM
    ResultadosConcreto AS r
JOIN
    Aditivos AS a ON r.aditivo_id = a.id
JOIN
    Productos AS p ON a.producto_id = p.id
GROUP BY
    p.nombre_producto, a.codigo, r.edad_dias
ORDER BY
    a.codigo, r.edad_dias;
\end{lstlisting}

Este diseño robusto de la base de datos es la piedra angular del proyecto, asegurando la calidad de los datos que alimentan el modelo matemático y las predicciones de la aplicación DIAPCE.

\subsubsection{Base de Datos Local del Cliente (SQLite)}

Continuando con la implementación, y para dar soporte a la funcionalidad de la aplicación en el dispositivo del usuario, se diseñó una segunda base de datos. A diferencia de la base de datos del servidor, que actúa como el repositorio central y definitivo de los datos experimentales, esta base de datos local, implementada en SQLite, tiene como propósito gestionar el estado de la aplicación, almacenar datos específicos del usuario y agilizar las consultas para las predicciones.

\paragraph{Esquema y Relaciones.}
El motor SQLite es gestionado en Flutter a través del paquete `sqflite`. El esquema se diseñó para incluir no solo las tablas de catálogo y los resultados experimentales (cargados localmente desde un archivo CSV para acceso rápido), sino también las tablas que manejan la lógica de la aplicación, como `users`, `projects` y `mixtures`.

Las relaciones entre estas tablas, como la que vincula un proyecto a un usuario, se establecen mediante claves foráneas con políticas de borrado en cascada para mantener la integridad de los datos del usuario.

\begin{figure}[H]
    \centering
    % --- INSTRUCCIÓN ---
    % Reemplaza la siguiente línea con la ruta a tu imagen del esquema completo de la BD.
    \includegraphics[width=1\textwidth]{images/diagrama_db_completo.jpg}
    \caption{Diagrama Entidad-Relación de la base de datos local completa, incluyendo las tablas de la aplicación y los datos experimentales.}
    \label{fig:sqlite_schema}
\end{figure}

El siguiente código DBML detalla la estructura completa de la base de datos local, mostrando todas las entidades y sus interconexiones.

\begin{lstlisting}[style=dbmlstyle, caption=Esquema completo de la base de datos local en DBML.]
Table users {
  id integer [pk, increment]
  email string [unique, not null]
  password string [not null]
}

Table projects {
  id integer [pk, increment]
  user_id integer [not null, ref: > users.id]
  project_name string [not null]
  selected_date string
  selected_image_path string
  creator_name string
  resistance_target integer [not null]  
  temperature integer [not null]        
  humidity integer [not null]           
  relacion_ac float [not null]          
  work_type string                      
  aditivo_id integer [ref: > Aditivos.id]
  resistencia_predicha_7d float
  resistencia_predicha_14d float
  resistencia_predicha_28d float
  created_at string
}

Table TiposAditivo {
  id integer [pk, increment]
  nombre string [unique, not null]
}

Table Productos {
  id integer [pk, increment]
  nombre_producto string [unique, not null]
  marca string
}

Table Aditivos {
  id integer [pk, increment]
  codigo string [unique, not null]         
  porcentaje_aplicado string [not null]    
  tipo_aditivo_id integer [not null, ref: > TiposAditivo.id]
  producto_id integer [not null, ref: > Productos.id]
}

Table ResultadosConcreto {
  id integer [pk, increment]
  temperatura integer [not null]           
  humedad integer [not null]               
  relacion_ac float [not null]             
  edad_dias integer [not null]             
  resistencia_mpa float [not null]
  aditivo_id integer [not null, ref: > Aditivos.id]
  
  indexes {
    (temperatura, humedad, relacion_ac, aditivo_id, edad_dias) [name: 'idx_busqueda_resistencia']
  }
}

Table mixtures {
  id integer [pk, increment]
  project_id integer [ref: > projects.id]
  name string
  description string
  cemento_kg float
  agua_litros float
  arena_kg float
  grava_kg float
  aditivo_porcentaje float
  created_at string
}
\end{lstlisting}

\paragraph{Flujo de Datos y Consultas para Predicciones.}
Al inicializar, la aplicación carga los datos desde un archivo CSV a la tabla `ResultadosConcreto`. Esta copia local permite que DIAPCE calcule las predicciones de resistencia de manera instantánea, sin depender de una conexión constante al servidor. Se creó un índice compuesto en la tabla `ResultadosConcreto` para optimizar la velocidad de estas consultas, que son el núcleo de la funcionalidad de la aplicación.

\begin{figure}[H]
    \centering
    % --- INSTRUCCIÓN ---
    % Reemplaza la siguiente línea con la ruta a tu imagen que ilustre el flujo de datos.
    \includegraphics[width=0.8\textwidth]{images/flujo_prediccion.jpg}
    \caption{Flujo de datos para el cálculo de predicciones en la base de datos local.}
    \label{fig:prediction_flow}
\end{figure}

La consulta fundamental que la aplicación ejecuta para obtener los promedios de resistencia a los 7, 14 y 28 días se basa en los parámetros seleccionados por el usuario y se muestra a continuación.

\begin{lstlisting}[style=sqlstyle, caption=Consulta clave para calcular promedios de resistencia.]
SELECT
    edad_dias,
    ROUND(AVG(resistencia_mpa), 2) AS resistencia_promedio,
    COUNT(*) AS num_muestras
FROM 
    resultados_concreto
WHERE 
    temperatura = ? AND humedad = ? AND relacion_ac = ? AND aditivo_id = ?
GROUP BY 
    edad_dias
ORDER BY 
    edad_dias ASC;
\end{lstlisting}

ya teniendo la informacion y los datos claros se muestran las intefaces
% Asegúrate de tener estos paquetes en el preámbulo de tu documento:
% \usepackage{graphicx}
% \usepackage{float}

\subsection{Vistas de la aplicación}

\subsubsection{Inicio de sesión (\texttt{lib/view/main\_login.dart})}
Función:
\begin{itemize}
  \item Autenticación local de usuarios (email + contraseña).
  \item Validación básica de campos y alternancia de visibilidad de contraseña.
  \item Navegación al panel principal (Hall) tras login exitoso.
\end{itemize}

\begin{figure}[H]
  \centering
  \includegraphics[width=0.5\linewidth]{images/main_login.jpg}
  \caption{Pantalla de inicio de sesión de la aplicación.}
  \label{fig:view-login}
\end{figure}

\subsubsection{Crear cuenta (\texttt{lib/view/create\_count.dart})}
Función:
\begin{itemize}
  \item Registro de nuevos usuarios con confirmación de contraseña.
  \item Validaciones de formato y coincidencia.
  \item Acceso a login tras registro.
\end{itemize}

\begin{figure}[H]
  \centering
  \includegraphics[width=0.5\linewidth]{images/create_count.jpg}
  \caption{Pantalla de registro de nuevos usuarios.}
  \label{fig:view-registro}
\end{figure}

\subsubsection{Panel de proyectos (Hall) (\texttt{lib/view/hall.dart})}
Función:
\begin{itemize}
  \item Listado en cuadrícula de los proyectos del usuario.
  \item Acciones rápidas: crear nuevo proyecto, abrir detalle, eliminar con confirmación.
  \item Carga reactiva desde el servicio de proyectos.
\end{itemize}

\begin{figure}[H]
  \centering
  \includegraphics[width=0.5\linewidth]{images/hall_vacio.jpg}
  \caption{Panel principal (Hall) que muestra los proyectos del usuario (vacio).}
  \label{fig:view-hall}
\end{figure}
\begin{figure}[H]
  \centering
  \includegraphics[width=0.5\linewidth]{images/hall_lleno.jpg}
  \caption{Panel principal (Hall) que muestra los proyectos del usuario (lleno).}
  \label{fig:view-hall}
\end{figure}

\subsubsection{Crear proyecto (\texttt{lib/view/create\_proyect\_screen.dart})}
Función:
\begin{itemize}
  \item Flujo de decisión en cascada basado en datos: Temperatura $\rightarrow$ Humedad $\rightarrow$ Relación a/c $\rightarrow$ Aditivo (con códigos \texttt{P0}–\texttt{PP3}).
  \item Campo de \textit{resistencia objetivo} con decimales, validado en rango [24.00, 57.00] MPa.
  \item Cálculo de predicciones (7/14/28 días) a partir del dataset experimental antes de continuar.
\end{itemize}

\begin{figure}[H]
  \centering
  \includegraphics[width=0.5\linewidth]{images/create_proyect_screen.jpg}
  \caption{Pantalla de creación de un nuevo proyecto con el flujo en cascada.}
  \label{fig:view-crear-proyecto}
\end{figure}
\begin{figure}[H]
  \centering
  \includegraphics[width=0.5\linewidth]{images/create_proyect_screen_2.jpg}
  \caption{Pantalla de creación de un nuevo proyecto con el flujo en cascada segunda parte.}
  \label{fig:view-crear-proyecto}
\end{figure}

\subsubsection{Detalle de proyecto (\texttt{lib/view/ViewExistingProjectScreen.dart})}
Función:
\begin{itemize}
  \item Visualización integral del proyecto: parámetros seleccionados, predicciones 7/14/28 días.
  \item Gráficos: composición de mezcla (pastel) y evolución de resistencia (línea).
  \item Generación automática de mezcla si falta, y opción de guardar/actualizar.
\end{itemize}

\begin{figure}[H]
  \centering
  \includegraphics[width=0.5\linewidth]{images/ViewExistingProjectScreen.jpg}
  \caption{Detalle de un proyecto existente con gráficos de predicción.}
  \label{fig:view-detalle-proyecto}
\end{figure}
\begin{figure}[H]
  \centering
  \includegraphics[width=0.5\linewidth]{images/ViewExistingProjectScreen_2.jpg}
  \caption{Detalle de un proyecto existente con gráficos de predicción prt 2.}
  \label{fig:view-detalle-proyecto}
\end{figure}
\begin{figure}[H]
  \centering
  \includegraphics[width=0.5\linewidth]{images/ViewExistingProjectScreen_3.jpg}
  \caption{Detalle de un proyecto existente con gráficos de predicción prt 3.}
  \label{fig:view-detalle-proyecto}
\end{figure}

\subsubsection{Vista de propiedades (legado, opcional) (\texttt{lib/view/view\_properties.dart})}
Función:
\begin{itemize}
  \item Pantalla histórica de propiedades; actualmente desactivada/comentada.
  \item Útil como referencia de diseño o comparación con la vista activa de detalle.
\end{itemize}

\begin{figure}[H]
  \centering
  \includegraphics[width=0.5\linewidth]{images/view_properties.jpg}
  \caption{Vista de propiedades (legado).}
  \label{fig:view-propiedades-legado}
\end{figure}
como inicio 
% docs/discusion_datos.tex
\section{Discusión de datos}

Esta sección analiza los datos y resultados generados con DIAPCE en el contexto del Meta (Colombia), con énfasis en uso académico en laboratorios de materiales y semilleros de investigación de la UCC. La discusión se centra en la calidad y cobertura del conjunto de datos  del sistema para dosificación base y clasificación de aditivos (P0–PP3) con proyecciones y (iii) la comparación frente a prácticas tradicionales de diseño según ACI 211.1.

\subsection{Conjunto de datos y preprocesamiento}
El conjunto de datos registra parámetros ambientales y de diseño (p. ej., resistencia objetivo, temperatura, humedad, condición de exposición) y variables de mezcla y desempeño. Para su uso en DIAPCE se aplican pasos mínimos de preprocesamiento: depuración de registros incompletos, homogeneización de unidades, detección de valores atípicos y codificación de la etiqueta de aditivo (P0–PP3). Se recomienda documentar las decisiones de limpieza empleadas en cada cohorte de ensayos para garantizar reproducibilidad.

\begin{table}[H]
  \centering
  \caption{Cobertura de rangos del conjunto de datos (ilustrativo)}
  \begin{tabular}{lccc}
    \hline
    Variable & Mín & Máx & Observaciones \\
    \hline
    Resistencia objetivo (MPa) & — & — & Ver distribución por percentiles \\
    Temperatura (°C) & — & — & \% de casos en tramos de 5 °C \\
    Humedad relativa (\%) & — & — & \% de casos en tramos de 10 \% \\
    Condición de exposición & — & — & Clases balanceadas/desbalanceadas \\
    \hline
  \end{tabular}
\end{table}

\subsection{Metodología de validación}
Para evaluar la validez de las recomendaciones, se propone una validación cruzada o partición entrenamiento/validación con un conjunto de retención para verificación externa. Las métricas consideradas son:
\begin{itemize}
  \item Clasificación de aditivos (P0–PP3): exactitud, precisión, exhaustividad y F1 macro; matriz de confusión por clase.
  \item Proyección de resistencia (7/14/28 días): MAE, RMSE y MAPE por horizonte.
  \item Coherencia normativa: cumplimiento de restricciones y rangos indicativos según ACI 211.1.
  \item Trazabilidad: \% de recomendaciones con insumos/resultados/versionado correctamente almacenados en SQLite.
\end{itemize}

\begin{table}[H]
  \centering
  \caption{Métricas de clasificación de aditivos (ilustrativo)}
  \begin{tabular}{lcccc}
    \hline
    Clase & Precisión & Exhaustividad & F1 & Soporte \\
    \hline
    P0 & — & — & — & — \\
    P1 & — & — & — & — \\
    P2 & — & — & — & — \\
    P3 & — & — & — & — \\
    Macro-promedio & — & — & — & — \\
    \hline
  \end{tabular}
\end{table}

\begin{table}[H]
  \centering
  \caption{Error de proyección de resistencia (7/14/28 días, ilustrativo)}
  \begin{tabular}{lccc}
    \hline
    Horizonte & MAE (MPa) & RMSE (MPa) & MAPE (\%) \\
    \hline
    7 días & — & — & — \\
    14 días & — & — & — \\
    28 días & — & — & — \\
    \hline
  \end{tabular}
\end{table}

\subsection{Resultados y análisis}
Se espera que DIAPCE proporcione: (i) sugerencias de dosificación base acordes con condiciones de entrada; (ii) clasificación P0–PP3 consistente con patrones observados; y (iii) proyecciones de resistencia razonables a 7/14/28 días. La interpretación debe considerar la cobertura de rangos: el desempeño es más confiable en regiones densamente muestreadas del conjunto de datos. Se recomienda reportar errores estratificados por tramos de temperatura, humedad y nivel de resistencia objetivo.



\subsection{Sostenibilidad e impacto}
Aunque DIAPCE no implementa optimización automática de proporciones, su uso sistemático puede incentivar decisiones de dosificación más consistentes y selección adecuada de aditivos. Cuando se disponga de mediciones, puede evaluarse el impacto potencial en: (i) contenido de cemento frente a una línea base docente, y (ii) variación de resistencia a 28 días. Estos análisis deben interpretarse dentro de los márgenes normativos y del contexto de formación.

\subsection{Amenazas a la validez y limitaciones}
\begin{itemize}
  \item Sesgo de datos: el conjunto es regional; la generalización fuera del Meta debe justificarse.
  \item Medición: variabilidad experimental en laboratorio puede introducir ruido en las proyecciones a 7/14/28 días.
  \item Cobertura: desempeño reducido en combinaciones poco representadas (p. ej., extremos de temperatura/humedad).
  \item Alcance del sistema: la herramienta asiste decisiones; no sustituye criterio profesional ni verificación normativa.
\end{itemize}

\subsection{Reproducibilidad}
DIAPCE almacena localmente entradas y resultados (SQLite), lo que permite rastrear cada recomendación respecto a sus insumos y la versión del modelo. Se sugiere anexar en el repositorio tablas exportables (CSV) con métricas por cohorte de ensayos y capturas del flujo de uso para respaldar la discusión.


% docs/conclusiones.tex
\section{Conclusiones}

Con base en los objetivos y el alcance real de DIAPCE, se concluye lo siguiente:

\begin{itemize}
  \item \textbf{Implementación de la herramienta:} Se desarrolló DIAPCE como una aplicación multiplataforma en Flutter con almacenamiento local mediante SQLite. La herramienta asiste la toma de decisiones en el diseño de mezclas: registra parámetros de entrada, sugiere una dosificación base y clasifica el uso de aditivos (P0--PP3), entregando proyecciones de resistencia en días en concordancia con ACI~211.1 y normas relacionadas.
  \item \textbf{Enfoque académico y trazabilidad:} Se consolidó un flujo de trabajo orientado a laboratorios de materiales y semilleros de investigación de la UCC, enfatizando trazabilidad (registro de insumos, resultados y versión del modelo) para favorecer reproducibilidad, evaluación docente y comparación entre ensayos.
  \item \textbf{Validez técnica en contexto formativo:} La validación propuesta --mediante métricas de clasificación (P0--PP3) y error de proyección (7/14/28)-- ofrece un marco objetivo para juzgar el desempeño del sistema. Los resultados esperados son consistentes con la normativa y adecuados para uso formativo, particularmente en los rangos del conjunto de datos con mayor cobertura. Los valores concretos deben reportarse en la sección de resultados correspondiente.
  \item \textbf{Adopción y operación:} El diseño offline-first y multiplataforma reduce fricciones de despliegue en entornos con conectividad limitada y equipos heterogéneos, facilitando su adopción en aula y laboratorio.
  \item \textbf{Contribución y límites:} DIAPCE no sustituye el criterio profesional ni el cumplimiento normativo; lo complementa con evidencia sistematizada y accesible. El uso disciplinado de la herramienta promueve decisiones basadas en datos y una cultura de registro. Entre las limitaciones se encuentran la cobertura regional del dataset y la ausencia de un módulo de optimización automática o de sincronización cliente-servidor, lo que a su vez define líneas claras de trabajo futuro.
\end{itemize}
    
\end{itemize}

% docs/recomendaciones.tex
\section{Recomendaciones}

A partir del desarrollo y uso previsto de DIAPCE, se proponen las siguientes recomendaciones:

\begin{itemize}
  \item \textbf{Ampliar y balancear el conjunto de datos:} Incorporar más ensayos regionales, equilibrar clases de aditivo (P0--PP3) y cubrir mejor los rangos de temperatura, humedad y resistencias objetivo. Documentar sistemáticamente los procesos de limpieza y criterios de exclusión.
  \item \textbf{Evaluación continua:} Establecer un flujo de métricas reproducible (exactitud/precisión/recall/F1 para P0--PP3; MAE/RMSE/MAPE para 7/14/28; coherencia ACI) con exportación a CSV/LaTeX para informes periódicos en laboratorio.
  \item \textbf{Módulos de valor agregado:} Añadir, como evolución, un análisis comparativo de costo y una estimación de huella de carbono (p. ej., a partir del contenido de cemento y factores de emisión), respetando límites normativos.
  \item \textbf{Validación en campo:} Ejecutar estudios de caso en obras reales en el Meta para contrastar proyecciones con resistencias ensayadas y retroalimentar la calibración del modelo.
  \item \textbf{Escalabilidad opcional:} Evaluar a futuro una capa de sincronización/servicios (backend) para consolidar datos entre dispositivos, gestionar usuarios y versionado central de modelos, manteniendo la operación offline como principio.
  \item \textbf{Aprendizaje automático (opcional):} Explorar modelos supervisados para mejorar la clasificación de aditivos y las proyecciones de resistencia, con validación cruzada, control de sobreajuste e interpretabilidad.
  \item \textbf{Usabilidad y docencia:} Incorporar exportación a PDF/CSV de proyectos y resultados, plantillas de prácticas de laboratorio, tutoriales integrados y consideraciones de accesibilidad, reforzando su utilidad educativa.
  \item \textbf{Generalización regional:} Adaptar el sistema para otras regiones de Colombia mediante parametrización de condiciones ambientales/datos locales y verificación normativa contextual.
  \item \textbf{Gobernanza de datos:} Mantener versionado de datasets y modelos, control de calidad de datos y lineamientos éticos para el uso de información experimental y recomendaciones.
\end{itemize}


\printbibliography

\newpage
\appendix
\section{Anexos}

\begin{table}[h!]
    \centering
    \caption{Cronograma de actividades}
    \label{tab:cronograma}
    \small % Hacemos la fuente un poco más pequeña para que quepa mejor
    \begin{tabular}{|p{3.2cm}|c|c|c|c|c|c|c|c|c|c|c|c|}
        \hline
        \textbf{Descripción de actividades} & \multicolumn{12}{c|}{\textbf{Meses}} \\
        \cline{2-13}
         & \textbf{1} & \textbf{2} & \textbf{3} & \textbf{4} & \textbf{5} & \textbf{6} & \textbf{7} & \textbf{8} & \textbf{9} & \textbf{10} & \textbf{11} & \textbf{12} \\
        \hline
        Análisis y levantamiento de requisitos & $\blacksquare$ & $\blacksquare$ & & & & & & & & & & \\
        \hline
        Diseño del modelo de datos & & & $\blacksquare$ & $\blacksquare$ & & & & & & & & \\
        \hline
        Implementación del modelo matemático & & & & $\blacksquare$ & $\blacksquare$ & $\blacksquare$ & & & & & & \\
        \hline
        Arquitectura del software & & & & & & $\blacksquare$ & $\blacksquare$ & $\blacksquare$ & & & & \\
        \hline
        Pruebas y evaluación & & & & & & & & $\blacksquare$ & $\blacksquare$ & $\blacksquare$ & $\blacksquare$ & $\blacksquare$ \\
        \hline
    \end{tabular}
\end{table}

\vspace{1cm}

\begin{table}[h!]
    \centering
    \caption{Presupuesto}
    \label{tab:presupuesto}
    % Usamos columnas p{} para permitir el ajuste de texto y controlar el ancho
    \begin{tabular}{|p{2.5cm}|p{3cm}|>{\centering\arraybackslash}p{1.5cm}|>{\raggedleft\arraybackslash}p{2.5cm}|>{\raggedleft\arraybackslash}p{2.5cm}|}
        \hline
        \textbf{RUBROS} & \textbf{Justificación} & \textbf{Cantidad}  & \textbf{Valor unitario} & \textbf{Valor total} \\
        \hline
        \textbf{Talento humano} &  \\
        & Jóvenes investigadores & 1 & 4.000.000 COP & 4.000.000 COP \\
        & Asesor & & & \\
        \hline
        \textbf{Equipos y software} & Computador portátil & 1 & 3.000.000 COP & 3.000.000 COP \\
        \hline
        \textbf{Capacitación} & Curso de Flutter & 1 & 600.000 COP & 600.000 COP \\
        \hline
        \textbf{Otros (Imprevistos)} & Posible alojamiento del backend & 1 & 200.000 COP & 200.000 COP \\
        \hline
        \multicolumn{4}{|r|}{\textbf{Valor total}} & \textbf{7.800.000 COP} \\
        \hline
    \end{tabular}
\end{table}

\section{Bibliografía}

\begin{thebibliography}{99}

\bibitem{lucidchart2024}
Lucidchart. (2024). Cómo utilizar el desarrollo basado en funciones. Recuperado de \url{https://www.lucidchart.com/blog/es/utilizar-el-desarrollo-basado-en-funciones}

\bibitem{tashildar2020}
Tashildar, A., Shah, N., Gala, R., Giri, T., \& Chavhan, P. (2020).
Application development using Flutter.
International Research Journal of Modernization in Engineering Technology and Science, 2(8), 1262--1266.

\bibitem{meetupBogota}
Meetup. (s.f.). Flutter Bogotá. Meetup.

\bibitem{devtopFlutterCol}
Devtop. (s.f.). Servicios de desarrollo Flutter en Colombia. Devtop.

\bibitem{achilles2025}
Achilles. (s.f.). El sector de la construcción: Retos y desafíos para 2025.

\bibitem{cincodias2024a}
CincoDías. (2024a). De salarios premium a faltar 700{,}000 trabajadores: Los jóvenes huyen de la construcción.

\bibitem{cincodias2024b}
CincoDías. (2024b). Descarbonizar la construcción para salvar el planeta.

\bibitem{tomasNoticias}
Tomas Noticias. (s.f.). Cuáles son los desafíos del sector de la construcción en Colombia.

\bibitem{jimenez2020}
Jiménez, L., \& Alvarado, M. (2020).
Modelos matemáticos para la optimización de mezclas de concreto: una revisión.
Revista de Ingeniería Civil, 15(2), 45--58.

\bibitem{garcia2021}
García, P., \& Torres, R. (2021).
Integración de materiales cementantes alternativos en el diseño de mezclas de concreto.
Journal of Sustainable Construction, 8(3), 112--125.

\bibitem{torres2023}
Torres, L., \& Jiménez, F. (2023).
Aplicación de redes neuronales en la predicción de propiedades del concreto.
Inteligencia Artificial en la Construcción, 5(2), 78--90.

\bibitem{concrebrain}
ConcreBrain. (s.f.). Software desarrollado para predecir la resistencia del concreto basándose en parámetros personalizados mediante algoritmos de redes neuronales.

\bibitem{concreteMixDesign}
Concrete Mix Design. (s.f.). Software comercial para el diseño de mezclas de concreto utilizando algoritmos básicos de cálculo.

\bibitem{concretelab}
Concretelab. (s.f.). Herramienta que facilita el diseño y análisis de mezclas de concreto.

\bibitem{ASTM_C39}
ASTM C39/C39M. (s.f.). Standard Test Method for Compressive Strength of Cylindrical Concrete Specimens.

\bibitem{ASTM_D1238}
ASTM D1238. (s.f.). Standard Test Method for Melt Flow Rates of Thermoplastics by Extrusion Plastometer.

\bibitem{ASTM_D4404}
ASTM D4404-18. (2018). Standard Test Method for Determination of Pore Volume and Pore Volume Distribution of Soil and Rock by Mercury Intrusion Porosimetry.

\bibitem{ASTM_E18}
ASTM E18-20. (2020). Standard Test Methods for Rockwell Hardness of Metallic Materials.

\bibitem{dartdev}
Dart.dev. (s.f.). Dart Programming Language. Recuperado de \url{https://dart.dev}

\bibitem{flutterdev}
Flutter.dev. (s.f.). Flutter — Build apps for any screen. Recuperado de \url{https://flutter.dev}

\bibitem{githubdocs}
GitHub Docs. (s.f.). GitHub Documentation. Recuperado de \url{https://docs.github.com}

\bibitem{ISO1183}
ISO 1183-1:2019. (2019). Plastics — Methods for determining the density of non-cellular plastics.

\bibitem{typescript}
TypeScriptlang.org. (s.f.). TypeScript: Typed JavaScript at Any Scale. Recuperado de \url{https://www.typescriptlang.org}

\bibitem{vscodedocs}
Visual Studio Code Docs. (s.f.). Visual Studio Code Documentation. Recuperado de \url{https://code.visualstudio.com/docs}

\bibitem{fddsite}
Feature Driven Development. (s.f.). Recuperado de \url{https://featuredrivendevelopment.com}

\bibitem{palmer2002}
Palmer, S. R., \& Felsing, J. M. (2002).
A Practical Guide to Feature-Driven Development.
Prentice Hall.

% --- Referencias técnicas añadidas (DIAPCE) ---
\bibitem{ACI211}
ACI Committee 211. (2016).
ACI 211.1-16: Standard Practice for Selecting Proportions for Normal, Heavyweight, and Mass Concrete.
American Concrete Institute.

\bibitem{ASTM_C494}
ASTM C494/C494M. (2022).
Standard Specification for Chemical Admixtures for Concrete.
ASTM International.

\bibitem{ASTM_C192}
ASTM C192/C192M. (2021).
Standard Practice for Making and Curing Concrete Test Specimens in the Laboratory.
ASTM International.

\bibitem{sqliteDocs}
SQLite Consortium. (s.f.). SQLite Documentation. Recuperado de \url{https://www.sqlite.org/docs.html}

\bibitem{sqflite}
Tekartik. (s.f.). sqflite — SQLite plugin for Flutter. Recuperado de \url{https://pub.dev/packages/sqflite}

\end{thebibliography}

\end{document}